本部分首先介绍语言建模与机器学习的基础。我们从基本原理出发，说明如何用词元和嵌入表示语言；随后讨论注意力机制。这一概念既是 Transformer 架构的基础，也是大语言模型取得成功的关键。我们还将回顾两种重要的机器学习技术：监督学习与强化学习。下一章转向具体实践，介绍典型训练流程中的各个环节，从预训练、微调直到推理，并讨论常用的基准测试。本部分最后一章理论内容着眼于大语言模型中的“大”：我们讨论模型性能如何随参数量、训练数据和计算量扩展，以及如何在控制成本的同时利用规模优势，构建高效模型。本部分以一个实践案例收尾。你将从零开始创建一个微型大语言模型，亲身了解这些技术如何结合起来。
